\subsection{碰撞阶段、能量损失与摩擦生热}

碰撞过程持续时间很短，通常先用动量守恒联系碰撞前后，再用恢复系数或能量条件描述材料性质。
完全弹性碰撞同时守恒动量和动能；完全非弹性碰撞只表示碰后没有相对速度，并不意味着系统能量
消失。对质量 $m_1,m_2$ 的一维完全非弹性碰撞，动能损失为
\begin{equation*}
 \Delta K_{\rm loss}=\frac12\mu(u_1-u_2)^2,
 \qquad \mu=\frac{m_1m_2}{m_1+m_2},
\end{equation*}
可见损失只由碰前相对速度和约化质量决定。若碰撞后还有弹簧压缩、物块爬升或摆动，应把
“碰撞瞬间”和“碰后运动”分成两个阶段，分别使用适合的守恒关系。

滑动摩擦把宏观机械能转化为接触物体的内能。两个接触面相对滑动距离为 $s_{\rm rel}$ 时，产生的
热量为
\begin{equation*}
 Q=f_k s_{\rm rel}.
\end{equation*}
这里必须使用相对路程，而不是任一物体对地路程。若研究对象只取其中一个物体，摩擦功可以为正
也可以为负；若把两个物体及接触面一起纳入系统，摩擦内力总功等于机械能减少量的负值，而内能
增加量为正。区分“单个物体的摩擦功”和“系统的摩擦生热”，可以避免传送带、叠块等问题中常见的
符号错误。
